\clearpage
\section{Expressive Power of Knowledge Graph Embeddings}

Consider the following common relational patterns:
\begin{itemize}
  \item \textbf{Symmetry:} \(A\) is married to \(B\), and \(B\) is married to
    \(A\).
  \item \textbf{Inverse:} \(A\) is teacher of \(B\), and \(B\) is student of
    \(A\).
  \item \textbf{Composition:} \(A\) is son of \(B\), and \(C\) is sister of
    \(B\), so \(C\) is aunt of \(A\).
\end{itemize}

\subsection{Relational Patterns in TransE}

For each pattern above, determine whether TransE can model it perfectly with
\[
  \vect h+\vect\ell=\vect t
\]
for every relation. Explain each answer. A zero relation embedding is not
allowed because it makes related entities indistinguishable.

\begin{solution}
Write \(\vect m\), \(\vect t\), and \(\vect s\) for the embeddings of
\emph{married to}, \emph{teacher of}, and \emph{student of}, respectively.

\textbf{Symmetry cannot be modeled.}  For two distinct entities \(A\) and
\(B\), symmetry would require
\[
  \vect a+\vect m=\vect b,
  \qquad
  \vect b+\vect m=\vect a.
\]
Substituting the first equation into the second gives
\[
  \vect a+2\vect m=\vect a,
\]
so \(\vect m=\vect 0\).  The first equation then also gives
\(\vect a=\vect b\).  This is precisely the forbidden degenerate solution;
therefore TransE cannot represent a non-trivial symmetric relation perfectly.

\textbf{Inverse relations can be modeled.}  The two facts require
\[
  \vect a+\vect t=\vect b,
  \qquad
  \vect b+\vect s=\vect a.
\]
They are consistent exactly when
\[
  \vect s=-\vect t.
\]
For example, take
\[
  \vect a=\begin{bmatrix}1\\0\end{bmatrix},\qquad
  \vect b=\begin{bmatrix}0\\1\end{bmatrix},\qquad
  \vect t=\begin{bmatrix}-1\\1\end{bmatrix},\qquad
  \vect s=\begin{bmatrix}1\\-1\end{bmatrix}.
\]
Both entities have unit norm, both relations are nonzero, and the two
equations hold.

\textbf{Composition can be modeled.}  Let \(\vect r_{\rm son}\),
\(\vect r_{\rm sister}\), and \(\vect r_{\rm aunt}\) denote the three relation
embeddings.  The facts in the question require
\[
  \vect a+\vect r_{\rm son}=\vect b,
  \qquad
  \vect c+\vect r_{\rm sister}=\vect b,
  \qquad
  \vect c+\vect r_{\rm aunt}=\vect a.
\]
Consequently, the required composition rule is
\[
  \vect r_{\rm sister}
  =\vect r_{\rm aunt}+\vect r_{\rm son}.
\]
For a concrete unit-norm realization, choose
\[
  \vect a=\begin{bmatrix}1\\0\end{bmatrix},\qquad
  \vect b=\begin{bmatrix}0\\1\end{bmatrix},\qquad
  \vect c=\begin{bmatrix}-1\\0\end{bmatrix},
\]
and set
\[
  \vect r_{\rm son}=\begin{bmatrix}-1\\1\end{bmatrix},\qquad
  \vect r_{\rm sister}=\begin{bmatrix}1\\1\end{bmatrix},\qquad
  \vect r_{\rm aunt}=\begin{bmatrix}2\\0\end{bmatrix}.
\]
All three relations are nonzero and all three facts are satisfied exactly.
\end{solution}

\subsection{Relational Patterns in RotatE}

RotatE represents each entity as \(d\) two-dimensional points, equivalently a
vector in \(\R^{2d}\), and represents a relation by \(d\) rotation angles. For
each \(i=0,\ldots,d-1\), applying relation \(\ell\) rotates
\((h_{2i},h_{2i+1})\) clockwise by angle \(\ell_i\). Entity embeddings have
unit \(\ell_2\) norm.

Determine whether RotatE can model symmetry, inverse relations, and composition
perfectly, and justify each answer.

\begin{solution}
Let \(R_{\theta}\) denote clockwise rotation by angle \(\theta\).  Rotations
compose by adding their angles:
\[
  R_{\alpha}R_{\beta}=R_{\alpha+\beta},
\]
where angles are understood modulo \(2\pi\), coordinatewise when \(d>1\).

\textbf{Symmetry can be modeled.}  It requires
\[
  R_m\vect a=\vect b,
  \qquad
  R_m\vect b=\vect a.
\]
Thus applying the relation twice must give the identity.  A nonzero choice is
\(m=\pi\).  For example, with \(d=1\),
\[
  \vect a=\begin{bmatrix}1\\0\end{bmatrix},
  \qquad
  \vect b=\begin{bmatrix}-1\\0\end{bmatrix},
\]
a rotation by \(\pi\) swaps \(A\) and \(B\).

\textbf{Inverse relations can be modeled.}  If the teacher relation has
angle \(t\), choose the student relation to have angle
\[
  s\equiv -t\pmod{2\pi}.
\]
Then \(R_sR_t=I\).  For instance, \(t=\pi/2\) and \(s=-\pi/2\) are both
nonzero and map an entity to another unit vector and back.

\textbf{Composition can be modeled.}  The three facts give
\[
  R_{r_{\rm son}}\vect a=\vect b,
  \qquad
  R_{r_{\rm sister}}\vect c=\vect b,
  \qquad
  R_{r_{\rm aunt}}\vect c=\vect a.
\]
Therefore it suffices to choose
\[
  r_{\rm sister}
  \equiv r_{\rm aunt}+r_{\rm son}\pmod{2\pi}.
\]
For example, take \(d=1\), \(\vect c=(1,0)^\top\),
\(r_{\rm aunt}=\pi/3\), \(r_{\rm son}=\pi/2\), and
\(r_{\rm sister}=5\pi/6\).  Define \(\vect a=R_{\pi/3}\vect c\) and
\(\vect b=R_{\pi/2}\vect a\).  All entities have unit norm, all relation
angles are nonzero, and the composition holds exactly.
\end{solution}

\subsection{A Limitation of RotatE}

Give an example of a graph that RotatE cannot model perfectly. Determine
whether TransE can model the same graph. Assume relation embeddings cannot be
zero in either model.

\begin{solution}
Let
\[
  \mathcal E=\{a,b,c,d\},\qquad
  \mathcal L=\{r,s\},
\]
and let the only positive triples be
\[
  \mathcal S
  =\{(a,r,b),\ (a,r,c),\ (b,s,d)\}.
\]
In particular, \((c,s,d)\notin\mathcal S\).

Suppose RotatE had a perfect embedding, and write \(R_r\) and \(R_s\) for
the clockwise rotations associated with the two relations.  The first two
positive triples imply
\[
  R_r\vect a=\vect b,
  \qquad
  R_r\vect a=\vect c,
\]
and hence \(\vect b=\vect c\).  The third positive triple then implies
\[
  R_s\vect c
  =R_s\vect b
  =\vect d.
\]
Thus RotatE necessarily predicts the nonexistent triple \((c,s,d)\), a
contradiction.  This proof holds in every dimension and does not rely on
either relation being zero.

TransE cannot model this graph perfectly either.  Indeed,
\[
  \vect a+\vect r=\vect b,
  \qquad
  \vect a+\vect r=\vect c
\]
again forces \(\vect b=\vect c\).  From
\(\vect b+\vect s=\vect d\), it follows that
\[
  \vect c+\vect s=\vect d,
\]
which incorrectly predicts \((c,s,d)\).  Both models apply a fixed,
deterministic transformation to a head entity for each relation, so neither
can perfectly represent this one-to-many relation while keeping \(B\) and
\(C\) distinguishable.
\end{solution}
